\documentclass[11pt]{article}
\usepackage[a4paper,margin=29mm]{geometry}
\usepackage{amsmath,amssymb,amsthm,mathtools}
\usepackage[T1]{fontenc}
\usepackage{lmodern,microtype}
\usepackage[colorlinks=true,linkcolor=blue,citecolor=blue,urlcolor=blue]{hyperref}
\newtheorem{theorem}{Theorem}[section]
\newtheorem{lemma}[theorem]{Lemma}
\newtheorem{corollary}[theorem]{Corollary}
\newtheorem{proposition}[theorem]{Proposition}
\theoremstyle{remark}\newtheorem{remark}[theorem]{Remark}
\newcommand{\Q}{\mathbb Q}\newcommand{\R}{\mathbb R}
\newcommand{\PP}{\mathbb P}\newcommand{\Affine}{\mathbb A}
\title{Rational fibers of cubic maps and a sharp obstruction to scalar height bounds}
\author{}
\date{9 November 2026}
\begin{document}
\maketitle
\begin{abstract}
A rational function of two rational variables whose numerator and denominator
have total degree at most three has an infinite rational fiber in every
nonempty real open set.  The points on this fiber have unbounded height.
The assertion for smooth cubic fibers follows from the tangent map and
Mazur's torsion theorem.  Thus even a completely ineffective bound on the
height of a rational pair in terms of a single rational measurement requires
degree at least four, if the bound is to hold on a real open set.
This degree threshold is sharp: the rational fibers of $U^4+V^4$ are all
finite, by Faltings' theorem.  Both assertions extend to every number
field, using Merel's torsion theorem and simultaneous archimedean
approximation.  These assertions concern ambient rational
pairs and do not exclude stronger height bounds on special sets of elliptic
curve multiples.
\end{abstract}

\section{The height question}
For a rational pair $(u,v)$, write $(1:u:v)=(c:a:b)$ with coprime integers
$a,b,c$ and $c>0$, and put
\[
h_2(u,v)=\log\max\{|a|,|b|,c\}.
\]
For $q=a/b\in\Q$ in lowest terms with $b>0$, let
$h(q)=\log\max\{|a|,b\}$.  Only finitely many rational pairs have
$h_2$ at most any fixed real number.

\begin{theorem}\label{thm:main}
Let $A,B\in\Q[U,V]$ be relatively prime, with $B\ne0$,
$\max\{\deg A,\deg B\}\le3$, and $A/B$ nonconstant.
For every nonempty open subset $\Omega\subset\R^2$, there exists $c\in\Q$
such that
\[
\{(u,v)\in\Omega\cap\Q^2:B(u,v)\ne0,\ A(u,v)/B(u,v)=c\}
\]
is infinite.  Its $h_2$-heights are unbounded.
\end{theorem}

The open set may in particular impose any finite list of strict real
inequalities.  Thus positivity conditions on two real conjugates do not
remove the obstruction in this theorem.

\section{The tangent map on a plane cubic}
Let $C\subset\PP^2$ be a smooth plane cubic over $\Q$ possessing a rational
point.  Its tangent map $\tau:C\to C$ sends $P$ to the residual point of
intersection with the tangent line at $P$, with intersection multiplicities
included.  This is a morphism defined over $\Q$.  Let $H$ denote the
hyperplane divisor class on $C$.

\begin{lemma}\label{lem:torsion}
For $P\in C(\Q)$, put $D_P=H-3[P]\in\operatorname{Pic}^0(C)(\Q)$.
If $D_P$ is torsion, then
\begin{equation}\label{eq:iterate}
\tau^{39}(P)=\tau^3(P).
\end{equation}
The two morphisms $\tau^{39}$ and $\tau^3$ are distinct on every smooth
plane cubic over a field of characteristic zero.
\end{lemma}

\begin{proof}
Choose an origin $O\in C(\Q)$ and identify $C$ with its Jacobian $E$.
Write $H_0=H-3[O]\in E(\Q)$.  The tangent intersection identity
$2[P]+[\tau(P)]\sim H$ gives
\[
\tau(P)=H_0-2P,
\qquad
\tau^k(P)=(-2)^kP+\frac{1-(-2)^k}{3}H_0.
\]
The coefficient in this formula is an integer.  Since $D_P=H_0-3P$,
\[
\tau^{39}(P)-\tau^3(P)
 =\frac{(-2)^3-(-2)^{39}}3D_P
 =183251937960D_P.
\]
Mazur's theorem \cite[Theorem~(7$'$)]{Mazur} says that every rational
torsion point of an elliptic curve has order in
$\{1,2,\ldots,10,12\}$.  All these integers divide
\[
2520=\operatorname{lcm}(1,2,\ldots,10,12),
\qquad
183251937960=2520\cdot72719023.
\]
This proves \eqref{eq:iterate} for torsion $D_P$.
Finally $\tau$ is translation followed by multiplication by $-2$; hence
$\deg\tau^k=4^k$.  The degrees for $k=39$ and $k=3$ are unequal, so
the morphisms cannot coincide.
\end{proof}

\begin{lemma}\label{lem:near}
If $C/\Q$ is a smooth projective genus-one curve, $P\in C(\Q)$, and
$E=\operatorname{Pic}^0(C)$ has a rational point of infinite order, then
every real neighborhood of $P$ in $C(\R)$ contains infinitely many
rational points of $C$.
\end{lemma}

\begin{proof}
Use $P$ as origin and let $D\in E(\Q)$ have infinite order.
The compact group $E(\R)$ contains a sequence of distinct multiples
$[m_j]D$ converging to its identity.  To see this directly, compactness
gives a convergent subsequence of the distinct points $[n]D$; differences
of suitable pairs give distinct nonzero multiples converging to the
identity.  Their translates by $P$ give the required rational points.
\end{proof}

\section{Proof of the fiber theorem}
\begin{proof}[Proof of Theorem~\ref{thm:main}]
Put $f=A/B$.  Delete the zero set of $B$ and the critical locus $df=0$.
Because $f$ is nonconstant in characteristic zero, what remains is a
nonempty Zariski open subset $W\subset\Affine^2$.  Every nonempty real
open subset of $\R^2$ meets $W(\Q)$.

First suppose the projective closure of the generic fiber of $f$ is not
a smooth geometrically integral plane cubic.  Then there is a nonempty
open subset $T\subset\Affine^1$ such that the fibers over $T$ have this same
property and degree at most three.  Choose
$P\in\Omega\cap W(\Q)$ with $c=f(P)\in T(\Q)$.
The geometric component of $A-cB=0$ through $P$ is unique because $P$ is
a smooth point of the fiber.  It is fixed by Galois, and is therefore
defined over $\Q$.  It is either a line, a conic, or a singular integral
plane cubic.  Its normalization has genus zero, and the smooth rational
point $P$ lifts to a rational point on that normalization.  The
normalization is consequently isomorphic to $\PP^1_\Q$.
A rational parametrization gives infinitely many rational points
arbitrarily close to $P$; all sufficiently close points belong to
$\Omega$ and avoid $B=0$.

Now suppose the generic projective fiber is a smooth geometrically
integral plane cubic.  After deleting finitely many values of $f$,
the projective fibers are smooth plane cubics.  Tangent intersection in
these fibers defines a rational map $\tau$ on the total space, defined
over $\Q$ and preserving $f$.  Equivalently it defines a rational map
of the plane on a nonempty open subset, since the total space of the
affine pencil is the graph of $f$.
The rational maps $\tau^{39}$ and $\tau^3$ are unequal: their
restrictions to a generic smooth cubic fiber have degrees $4^{39}$ and
$4^3$, respectively.  Thus the locus on which both are defined and
different contains a nonempty Zariski open subset of the plane.
Choose $P\in\Omega\cap\Q^2$ in this open subset, in $W$, and above a
smooth cubic fiber $C$.
Lemma~\ref{lem:torsion} shows that $D_P\in\operatorname{Pic}^0(C)(\Q)$
has infinite order.  Lemma~\ref{lem:near} supplies infinitely many
rational points of $C$ arbitrarily close to $P$.  A small enough
neighborhood stays in the affine chart, in $\Omega$, and away from
$B=0$.  This proves the required infinitude in both cases.

The height assertion follows because a bounded-height set of rational
pairs is finite.
\end{proof}

\begin{corollary}\label{cor:height}
Under the hypotheses of Theorem~\ref{thm:main}, there is no finite-valued
function $\Phi:[0,\infty)\to\R$ such that
\[
h_2(u,v)\le\Phi\bigl(h(f(u,v))\bigr)
\]
for every $(u,v)\in\Omega\cap\Q^2$ with $B(u,v)\ne0$.
In particular, no such $f$ is finite-to-one on the rational points of
any nonempty real open subset of its domain.
\end{corollary}

\begin{proof}
The infinite fiber in Theorem~\ref{thm:main} has one fixed value of
$h(f(u,v))$ and unbounded $h_2(u,v)$.
\end{proof}

\section{Sharpness of the degree threshold}
\begin{proposition}\label{prop:quartic}
The polynomial $g(U,V)=U^4+V^4$ is finite-to-one on $\Q^2$.
Consequently there exists a finite-valued nondecreasing function
$\Phi:[0,\infty)\to[0,\infty)$ such that
\[
h_2(u,v)\le\Phi\bigl(h(u^4+v^4)\bigr)\qquad(u,v\in\Q).
\]
This assertion gives no effective formula for $\Phi$.
\end{proposition}

\begin{proof}
The fiber above zero consists only of $(0,0)$.  For $c\in\Q^\times$,
the projective curve
\[
U^4+V^4=cW^4
\]
is smooth over $\overline\Q$: its three partial derivatives vanish
simultaneously only when $U=V=W=0$, which is not a projective point.
A smooth plane quartic has genus three.  Its rational points are finite
by Faltings' theorem \cite{Faltings}, proving the first assertion.
For each real $T\ge0$ there are only finitely many $c\in\Q$ with
$h(c)\le T$.  The union of their rational fibers under $g$ is finite.
Take $\Phi(T)$ to be the maximum of $h_2$ over that finite union,
including zero in the maximum.  This also proves monotonicity.
\end{proof}

\begin{remark}
Finite fibers are substantially weaker than injectivity.
Proposition~\ref{prop:quartic} does not solve the polynomial-injection
problem for $\Q^2\to\Q$.  Poonen's theorem \cite{Poonen} obtains such
an injection under a form of the Bombieri--Lang conjecture.
\end{remark}

\section{Scope for rational height measurements}
Suppose a quadratic number field is written
$K=\Q(\sqrt d)$, and a height argument produces an element
$S=u+v\sqrt d$ with $u,v\in\Q$.  If the only restrictions on $u,v$
are strict inequalities in the real embeddings, those restrictions
describe a real open set.  Corollary~\ref{cor:height} rules out replacing
the two independent rational coordinates by any single rational
function with numerator and denominator of total degree at most three
while retaining even an ineffective uniform height bound.

If, however, $S$ is a prescribed sum of abscissas of integer multiples
of fixed elliptic points, its attainable coordinate pairs form a much
smaller arithmetic set.  Theorem~\ref{thm:main} makes no assertion that
its infinite fiber meets this set infinitely often.  A height theorem
on that special set may therefore hold, and must use the arithmetic
restriction rather than positivity alone.  Similarly, this result is
not a lower bound for the number of quantified variables in a
Diophantine definition of the integers.

\section{Extension to every number field}
Let $F$ be a number field and let
\[
F_\infty=\prod_{v\mid\infty}F_v,
\]
where one place is used for each conjugate pair of complex embeddings.
We give this product its usual real and complex topology and view $F$
diagonally in it.  A $\Q$-basis of $F$ identifies
$F\otimes_\Q\R$ with $F_\infty$ as a real vector space; the density of
$\Q^{[F:\Q]}$ proves that $F$ is dense in $F_\infty$.
The same argument shows that the $F$-points of every nonempty Zariski
open subset of $\Affine_F^2$ meet every nonempty open subset of
$F_\infty^2$.
For $u,v\in F$ the pair height is the absolute logarithmic projective
height $h(1:u:v)$.  Bounded-height subsets of $F^2$ are finite: for an
algebraic number of degree at most $[F:\Q]$ and bounded height, its
primitive minimal polynomial has bounded Mahler measure and hence
bounded integer coefficients.  There are only finitely many such
polynomials and their roots.  Apply this to $u$ and $v$, whose heights
are each at most $h(1:u:v)$.  This is the form of Northcott's
finiteness theorem used below.

\begin{theorem}\label{thm:numberfield}
Let $A,B\in F[U,V]$ be relatively prime, with $B\ne0$,
$\max\{\deg A,\deg B\}\le3$, and $A/B$ nonconstant.
For every nonempty open subset $\Omega\subset F_\infty^2$, there is
$c\in F$ such that
\[
\{(u,v)\in F^2\cap\Omega:B(u,v)\ne0,\ A(u,v)/B(u,v)=c\}
\]
is infinite and has unbounded absolute logarithmic height.
\end{theorem}

\begin{proof}
Merel's uniform torsion theorem \cite[Introduction, Corollary]{Merel}
provides an integer $M_F\ge1$ annihilating the torsion subgroup of
$E(F)$ for every elliptic curve $E/F$.  For instance, if $N$ is an
integer upper bound for the orders of all such torsion points, one
may take $M_F=\operatorname{lcm}(1,\ldots,N)$.
Write $3M_F=2^a u$ with $u$ odd, and choose a positive even integer
$h$ such that $2^h\equiv1\pmod u$.  Then
\[
M_F\ \bigm|\ \frac{(-2)^a-(-2)^{a+h}}3.
\]
On a smooth plane cubic, the tangent-map computation of
Lemma~\ref{lem:torsion} therefore shows
\[
D_P\text{ torsion}\quad\Longrightarrow\quad
\tau^{a+h}(P)=\tau^a(P).
\]
These maps have distinct degrees.  The unequal-iterate locus is thus
a nonempty Zariski open set in any pencil with smooth generic cubic
fiber.  Choose an $F$-point $P$ in this locus and in the prescribed
archimedean open set, avoiding the same finite exceptional set of
fibers as in the proof over $\Q$.  Then $D_P$ has infinite order.

The group $\prod_{v\mid\infty}E(F_v)$ is compact.  Its infinite cyclic
subgroup generated by $D_P$ has a sequence of distinct multiples
converging to the identity simultaneously at every archimedean
place.  Translation by $P$ gives infinitely many $F$-points of its
fiber inside $\Omega$, avoiding the denominator and the points at
infinity after shrinking the neighborhoods.

If the generic projective fiber is not a smooth geometrically
integral cubic, the component through a smooth $F$-point is a rational
curve, exactly as in the proof of Theorem~\ref{thm:main}.  Its
normalization is $\PP_F^1$, and density of $F$ in $F_\infty$ gives
infinitely many points simultaneously close to the chosen point.
Finally, bounded-height subsets of $F^2$ are finite by Northcott's
finiteness theorem, so the heights in the infinite fiber are unbounded.
\end{proof}

\begin{corollary}\label{cor:numberfieldsharp}
Over every number field $F$, degree four is the least possible degree
of a rational scalar measurement of two variables that has finite
fibers on all of $F^2$.  More precisely, no rational function of
numerator and denominator degrees at most three is finite-to-one on
the $F$-points of any nonempty archimedean open set, whereas for every
$\delta\in F^\times\setminus F^{\times4}$ the polynomial
\[
g_\delta(U,V)=U^4-\delta V^4
\]
is finite-to-one on $F^2$.
There is a finite-valued nondecreasing function $\Phi_F$ such that
$h(1:u:v)\le\Phi_F(h(g_\delta(u,v)))$ for all $u,v\in F$.
\end{corollary}

\begin{proof}
The lower bound follows from Theorem~\ref{thm:numberfield}.
The zero fiber of $g_\delta$ contains only $(0,0)$: a point with
$v\ne0$ would make $\delta=(u/v)^4$.
For $c\ne0$, the projective curve
$U^4-\delta V^4=cW^4$ is smooth of genus three.  Its $F$-points are
finite by Faltings' theorem.  Northcott's theorem now gives the
height function $\Phi_F$ by the same finite-union argument as in
Proposition~\ref{prop:quartic}.
Such a $\delta$ exists in every number field: any rational prime
unramified in $F$ has valuation one at every prime of $F$ above it,
and hence cannot be a fourth power in $F$.
\end{proof}

Neither the proof of Theorem~\ref{thm:numberfield} nor the sharpness
argument assigns a numerical height bound.  The assertions remain
separate from the problem of compressing a prescribed arithmetic
subset of elliptic multiples.

\begin{thebibliography}{9}
\raggedright
\bibitem{Mazur}
B.~Mazur, \emph{Modular curves and the Eisenstein ideal},
Publications Math\'ematiques de l'IH\'ES \textbf{47} (1977), 33--186,
Theorems~(7$'$) and~(8).
\href{https://doi.org/10.1007/BF02684339}{doi:10.1007/BF02684339}.

\bibitem{Faltings}
G.~Faltings, \emph{Endlichkeitss\"atze f\"ur abelsche Variet\"aten
\"uber Zahlk\"orpern}, Inventiones Mathematicae \textbf{73} (1983),
349--366, Satz~7, p.~365; erratum, \textbf{75} (1984), 381.
\href{https://doi.org/10.1007/BF01388432}{doi:10.1007/BF01388432}.

\bibitem{Merel}
L.~Merel, \emph{Bornes pour la torsion des courbes elliptiques sur les
corps de nombres}, Inventiones Mathematicae \textbf{124} (1996),
no.~1--3, 437--449.
\href{https://doi.org/10.1007/s002220050059}{doi:10.1007/s002220050059};
\href{https://perso.imj-prg.fr/wp-content/uploads/merel-pub/torsion.pdf}
{author's manuscript}.

\bibitem{Poonen}
B.~Poonen, \emph{Multivariable polynomial injections on rational numbers},
Acta Arithmetica \textbf{145} (2010), no.~2, 123--127.
\href{https://doi.org/10.4064/aa145-2-2}{doi:10.4064/aa145-2-2};
\href{https://math.mit.edu/~poonen/papers/QQ.pdf}{author's manuscript}.
\end{thebibliography}
\end{document}
